\section{Metodología}

Se implementaron en Python (módulo \texttt{src/}) las dos estrategias de
decisión sobre el mismo problema bi-objetivo, garantizando reproducibilidad
bit-a-bit mediante una semilla fija en cada llamada a \texttt{minimize()}.

\subsection{Estrategia A Posteriori: NSGA-II + post-selección}

Se aproxima el frente de Pareto completo con el algoritmo
NSGA-II~\cite{deb2002nsga2,blank2020pymoo}. Sobre el conjunto de \MetNSol{}
soluciones no dominadas resultantes, el tomador de decisiones aplica dos
criterios de selección \emph{a posteriori}:

\begin{itemize}[noitemsep]
  \item \textbf{Suma ponderada normalizada.} Se normaliza cada objetivo al
        rango $[0,1]$ usando los extremos del frente y se minimiza la
        combinación $0.7\,\tilde{f_1} + 0.3\,\tilde{f_2}$, otorgando 70\,\% de
        importancia al costo y 30\,\% al tiempo.
  \item \textbf{Lexicográfica.} Se prioriza estrictamente el Tiempo ($f_2$):
        se elige la solución del frente con menor $f_2$.
\end{itemize}

\subsection{Estrategia A Priori: GA mono-objetivo}

Aquí las preferencias se inyectan \emph{antes} de optimizar, transformando el
problema en mono-objetivo y resolviéndolo con un GA clásico:

\begin{itemize}[noitemsep]
  \item \textbf{Suma ponderada directa.} Con un \emph{escalamiento manual}
        aproximado a los rangos de cada objetivo, se minimiza
        \[
          g(\mathbf{x}) = 0.7 \cdot \frac{f_1(\mathbf{x})}{200}
                        + 0.3 \cdot \frac{f_2(\mathbf{x})}{50}.
        \]
  \item \textbf{Lexicográfica directa.} Se optimiza únicamente el objetivo
        prioritario, $f_2$, ignorando por completo a $f_1$ durante la búsqueda.
\end{itemize}

\subsection{Parámetros de ejecución}

La Tabla~\ref{tab:parametros} resume la configuración. Ambos algoritmos usan el
mismo tamaño de población y número de generaciones para una comparación justa.

\begin{table}[h]
  \centering
  \caption{Parámetros de los algoritmos evolutivos.}
  \label{tab:parametros}
  \begin{tabular}{lll}
    \toprule
    \textbf{Parámetro} & \textbf{A Posteriori (NSGA-II)} & \textbf{A Priori (GA)} \\
    \midrule
    Tamaño de población & \num{100} & \num{100} \\
    Generaciones        & \num{100} & \num{100} \\
    Semilla             & \MetSeed & \MetSeed \\
    Objetivos           & 2 ($f_1, f_2$) & 1 (escalar) \\
    Pesos $(w_1, w_2)$  & $(0.7,\ 0.3)$ & $(0.7,\ 0.3)$ \\
    Escalamiento        & normalización por rango & manual ($f_1/200,\ f_2/50$) \\
    \bottomrule
  \end{tabular}
\end{table}

\subsection{Métricas de calidad del frente}

Para caracterizar el frente NSGA-II se reportan: el \textbf{hipervolumen}
(volumen dominado respecto al punto de referencia $(\num{200}, \num{50})$),
indicador de convergencia y diversidad; y el \textbf{spacing} de
Schott~\cite{coello2007evolutionary}, que mide la uniformidad de la
distribución de soluciones a lo largo del frente.
